\section{Appendix - Experimental Details}\label{sec:appendix:experimental-details}

\subsection{Task and Data Generation}

Each task followed the A1~$\rightarrow$~B~$\rightarrow$~A2 continual learning schedule adapted from \citet{holton2025humans}. Tasks consisted of 6 plant cues with angular responses on a circular dial, with summer and winter targets differing by a fixed angular offset. Inputs were 12-dimensional one-hot vectors (6 cues $\times$ 2 seasons); outputs were four continuous values (cosine--sine encodings for summer and winter targets).

\paragraph{Experimental Conditions} Experiments varied: (1) 2 architectures (modular task-routed recurrent network; single recurrent baseline), (2) 5 initialization scales $\gamma \in \{0.001, 0.01, 0.1, 1.0, 2.0\}$, and (3) 3 task similarity conditions (Same, Near, Far). Each architecture $\times$ $\gamma$ $\times$ similarity condition contained 101--103 independent runs (103 Same, 101 Near, 101 Far schedules), yielding approximately 3050 trained networks.

\paragraph{Trials per Phase} Each schedule contained 60 summer and 60 winter trials per phase (120 total). With 100 epochs per phase, this yielded 6000 summer and 6000 winter updates for A1 and B. During A2, winter trials were evaluated but did not contribute gradients (held-out probe).

\subsection{Training}

All models used SGD (learning rate: 0.01, momentum: 0.0, weight decay: 0.0, no gradient clipping) with batch size 1, fixed trial order, and 100 epochs per phase. Optimizer state was preserved across phases.

Recurrent weights used PyTorch's default uniform initialization, rescaled by $\gamma \in \{0.001, 0.01, 0.1, 1.0, 2.0\}$, defining the rich-to-lazy regime. Recurrent biases were disabled; only the readout layer had bias. Random seed was fixed at 2024; individual networks received unique weights via sequential RNG sampling.

\subsection{Behavioral Measures}

\paragraph{Accuracy} Predicted cosine--sine outputs were converted to angle $\hat{\theta}$:
\[
\text{Accuracy} = 1 - \frac{\left|\mathrm{wrap}_{[-\pi,\pi]}(\hat{\theta}-\theta)\right|}{\pi}
\]
Learning curves used a centered rolling average (window = 25 trials) averaged across runs.

\paragraph{Transfer} Performance retention across A1~$\rightarrow$~B:
\[
\text{Transfer} = \overline{\text{Acc}}_{B,\text{winter},\text{first 6}} - \overline{\text{Acc}}_{A1,\text{winter},\text{last 6}}
\]

\paragraph{Interference} Measured from A2 winter responses via two-component von Mises mixture:
\[
p(\phi) = \pi_A \cdot \mathrm{vM}(\phi;\mu_A,\kappa) + \pi_B \cdot \mathrm{vM}(\phi;\mu_B,\kappa), \quad \text{Interference} = 1 - \pi_A
\]
where $\mu_A$, $\mu_B$ are task-A and task-B rules, and $\pi_A + \pi_B = 1$. Computed only for Near and Far conditions. All metrics reported as mean $\pm$ SEM.

\subsection{Architecture Details}

\begin{table}[h]
\centering
\small
\begin{tabular}{lccccc}
\toprule
Architecture & Modules & Units/module & Total hidden & Recurrent params & Readout params \\
\midrule
Modular & 2 & 25 & 50 & 1850 & 204 \\
Single & 1 & 50 & 50 & 3100 & 204 \\
\bottomrule
\end{tabular}
\end{table}

Architectures were state-matched (equal hidden dimensionality) but not parameter-matched (single network had ~1250 additional recurrent parameters). Recurrent layers had no biases; output layer bias was enabled. In the modular architecture, no recurrent weights were shared; each module had independent recurrent and input projections with task-routed input (inactive module received zero input).

All models used 1 recurrent layer with 2 unroll steps. The modular architecture had no inter-module recurrent connectivity (block-diagonal dynamics: $W_{\mathrm{rec}} = \mathrm{blockdiag}(W_A, W_B)$); modules interacted only through the shared readout. Ablations explored sparse inter-module coupling (densities 0.3, 0.5, 0.7, 1.0) with identical training parameters (not shown in main figures).
%%%%%%%%%%%%%%%%%%%%%%%%%%%%%%%%%%%%%%%%%%%%%%%%%%%%%%%%%%%%%%%%%%%%%%%%%%%%%%%%%%%%%%%%%%%%%%%%%%%%%%

\section{Appendix - Analysis Details}
\label{sec:appendix:analysis-details}

\paragraph{Angular reconstruction and accuracy.}
On each trial, the network produces a four-dimensional output corresponding to two cosine--sine pairs, one for the summer feature (\texttt{feature\_idx} = 0) and one for the winter feature (\texttt{feature\_idx} = 1), as described in the Methods section. For a given probed feature $f \in \{\text{summer}, \text{winter}\}$, we reconstruct the predicted angle
\[
\hat{\theta}_t^{(f)} = \operatorname{atan2}\!\big(\hat{s}_t^{(f)}, \hat{c}_t^{(f)}\big),
\]
where $(\hat{c}_t^{(f)}, \hat{s}_t^{(f)})$ are the cosine and sine outputs for that feature at trial $t$, and $\operatorname{atan2}$ denotes the two-argument arctangent. The target angle is denoted $\theta_t^{(f)}$, and we compute the circular error
\[
\Delta \theta_t^{(f)} = \operatorname{wrap}\big(\hat{\theta}_t^{(f)} - \theta_t^{(f)}\big),
\]
where $\operatorname{wrap}$ maps angles to the interval $[-\pi, \pi)$. Unless stated otherwise, \emph{accuracy} is one minus the mean normalized absolute circular error,
\[
\mathrm{Acc} = 1 - \frac{1}{T} \sum_{t=1}^T \frac{\lvert \Delta \theta_t^{(f_t)} \rvert}{\pi},
\]
computed over a specified set of trials (e.g., all winter trials in a given phase). Equivalently, the per-trial accuracy $\mathrm{Acc}_t = 1 - \lvert \Delta \theta_t^{(f_t)} \rvert / \pi$ is averaged over $T$ trials. The metric is bounded in $[0, 1]$, with higher values indicating better performance.

\paragraph{Transfer measure.}
Following Holton et al.~\cite{holton2025humans}, we quantify transfer from task A to task B using performance on winter trials. Let $\mathcal{W}_{A1}^{\text{end}}$ denote the set of the last six winter trials in phase $A1$, and let $\mathcal{W}_{B}^{\text{start}}$ denote the set of the first six winter trials in phase $B$. We define
\[
\mathrm{Acc}_{A1}^{\text{end}} = \frac{1}{\lvert \mathcal{W}_{A1}^{\text{end}} \rvert} \sum_{t \in \mathcal{W}_{A1}^{\text{end}}} \mathrm{Acc}_t,
\quad
\mathrm{Acc}_{B}^{\text{start}} = \frac{1}{\lvert \mathcal{W}_{B}^{\text{start}} \rvert} \sum_{t \in \mathcal{W}_{B}^{\text{start}}} \mathrm{Acc}_t,
\]
where $\mathrm{Acc}_t$ is the single-trial accuracy defined from $\Delta \theta_t^{(f_t)}$ as above. The \emph{transfer} metric is then
\[
\mathrm{Transfer} = \mathrm{Acc}_{B}^{\text{start}} - \mathrm{Acc}_{A1}^{\text{end}},
\]
so that positive values indicate that prior learning on task A facilitates early performance on task B. In practice, transfer values are typically negative (winter accuracy drops at the $A1 \rightarrow B$ boundary because the rule has changed); less-negative values therefore correspond to better transfer between tasks.

\paragraph{Mixture-based interference measure.}
To quantify interference at retest on task A, we analyze winter responses in phase $A2$ using a mixture model over the network's implicit \emph{rule angle} on each trial. For a given run and condition, and for every winter trial $i$ in phase $A2$, we define the rule-angle residual between the network's winter and summer angle predictions on matched plant cues,
\[
\phi_i \;=\; \operatorname{wrap}\!\big( \hat{\theta}_{\text{winter},i} - \hat{\theta}_{\text{summer},i} \big),
\]
where $(\hat{\theta}_{\text{summer},i}, \hat{\theta}_{\text{winter},i})$ are reconstructed from the network's two cosine--sine output channels (output indices 0--1 and 2--3, respectively). Let $\theta_A$ and $\theta_B$ denote the task-A and task-B rule angles (the angular offsets between summer and winter targets) stored from the schedule. We fit a two-component von~Mises mixture
\[
p(\phi) = \pi_A \, \mathrm{VM}(\phi \mid \theta_A, \kappa) + (1 - \pi_A) \, \mathrm{VM}(\phi \mid \theta_B, \kappa),
\]
where $\mathrm{VM}(\cdot \mid \mu, \kappa)$ is the von~Mises distribution with mean direction $\mu$ and concentration $\kappa$. For simplicity and identifiability, we use a shared concentration parameter $\kappa$ across components and fix the component means to $(\theta_A, \theta_B)$; the only free parameters are $(\pi_A, \kappa)$, with $\pi_B = 1 - \pi_A$.

We estimate $(\pi_A, \kappa)$ by maximizing the log-likelihood
\[
\mathcal{L}(\pi_A, \kappa) = \sum_{i=1}^N \log p(\phi_i)
\]
using an expectation--maximization (EM) procedure that alternates an analytic E-step for the membership weights with an L-BFGS-B M-step on $\kappa$, run to convergence (parameter-update tolerance $10^{-3}$). To avoid local optima, we sweep a deterministic grid of starting points: $\pi_A^{(0)} \in \{0.1, 0.2, \ldots, 0.9\}$ crossed with $\kappa^{(0)} \in \{1, 2.5, 5, 10, 15, 20\}$, giving 54 fits per run, and retain the fit with the highest log-likelihood. For each model run, we define the \emph{interference} index as
\[
\mathrm{Interference} = 1 - \pi_A^\star,
\]
where $\pi_A^\star$ is the estimated weight on the task-A component at the chosen optimum. Values closer to 1 indicate stronger shift of behavior toward the task-B rule at retest.

\paragraph{Uncertainty estimates.}
Uncertainty in behavioral and representational plots is reported by panel type. For the per-run summary statistics shown in the transfer, interference, effective-dimensionality, and principal-angle panels, we use bootstrap 95\% confidence intervals: simulation runs are resampled with replacement $B = 1{,}000$ times within each (architecture, $\gamma$, similarity) cell, the summary statistic of interest is recomputed on each bootstrap sample, and we report the point estimate together with the central 95\% interval (2.5th to 97.5th percentile) as error bar. For the smoothed accuracy and loss curves (Figures showing the $A1 \rightarrow B \rightarrow A2$ time course), we instead report the participant-level mean $\pm$ SEM at each trial, with both the mean and the SEM passed through a centred rolling-mean window of 25 trials. All hypothesis tests reported in the main text are descriptive and based on these distributions; we do not perform formal multiple-comparison corrections.

\paragraph{Effective dimensionality and principal angles.}
For each model run, architecture, similarity, and phase $\phi \in \{A1, B, A2\}$, we form a per-phase data matrix $H_\phi \in \mathbb{R}^{12 \times h}$ by stacking the final-time-step hidden states for the 12 canonical-sweep stimuli (6 task-A stimuli followed by 6 task-B stimuli) of that phase, where $h$ is the total hidden width. PCA is fit on $H_\phi$ alone, and the effective dimensionality at variance threshold $\tau$ is
\[
d_{\mathrm{eff}}(\tau) = \min\!\left\{ k \;:\; \frac{\sum_{j=1}^k \lambda_j}{\sum_{j=1}^{D} \lambda_j} \geq \tau \right\},
\]
where $\{\lambda_j\}$ are the eigenvalues of the centred covariance of $H_\phi$ sorted in decreasing order and $D = \min(12, h)$. We report $d_{\mathrm{eff}}(0.99)$ at $\phi = B$ in the corresponding figure panels. We verified that the qualitative patterns reported in the Results are robust to varying the threshold between 95\% and 99\%.

To compute principal angles between task-specific subspaces, we work from the post-$B$ matrix $H_B \in \mathbb{R}^{12 \times h}$ defined above and split it by task identity along the canonical stimulus sweep, yielding $H^{(A)} \in \mathbb{R}^{6 \times h}$ (first 6 rows: task A) and $H^{(B)} \in \mathbb{R}^{6 \times h}$ (last 6 rows: task B). We then fit PCA separately on $H^{(A)}$ and $H^{(B)}$ and retain the first two principal components of each, yielding two-dimensional subspaces $\mathcal{S}_A$ and $\mathcal{S}_B$ in the shared hidden space. Let $U_A, U_B \in \mathbb{R}^{h \times 2}$ be orthonormal bases of $\mathcal{S}_A, \mathcal{S}_B$, and let $\sigma_1 \geq \sigma_2$ be the singular values of $U_A^\top U_B$. We report the first (smallest) principal angle in degrees,
\[
\theta_1 = \arccos(\sigma_1),
\]
as our scalar summary. The first principal angle ranges from $0$ (when $\mathcal{S}_A$ and $\mathcal{S}_B$ share at least one direction) to $90^\circ$ (when $\mathcal{S}_A$ and $\mathcal{S}_B$ are mutually orthogonal in the ambient space); larger values therefore correspond to stronger separation of task-specific subspaces.

\paragraph{PCA visualizations.}
For the qualitative 3D PCA trajectories shown in the main figures, we fit a single PCA on hidden states concatenated across phases $A1$, $B$, and $A2$ for a representative run at each architecture and initialization regime, using the same canonical stimulus sweep as above. We then project the hidden states for each phase and task onto the first three principal components and plot the resulting trajectories as closed loops connecting stimuli in sweep order. Colors encode phase and task identity, as described in the figure captions. These visualizations are used solely for illustration and do not enter into any quantitative analysis.